% Un punto della nuvola e le triple che lo descrivono.
% Nomi reali di 3DOntCore (si veda grande_3DGraph.nt).
% Incluso da chapters/introduction.tex (sottosezione Vantaggi).
\begin{tikzpicture}[
  font=\footnotesize,
  iri/.style={draw, semithick, rounded corners=3pt, fill=black!6,
              inner sep=3pt, minimum height=5.5mm},
  cls/.style={draw, semithick, rounded corners=3pt, fill=black!14,
              inner sep=3pt, minimum height=5.5mm},
  lit/.style={draw, semithick, inner sep=3pt, minimum height=5.5mm,
              minimum width=20mm},
  arc/.style={-{Stealth[length=4.5pt]}, semithick},
  side/.style={font=\scriptsize, inner sep=1pt, fill=white},
]
  % catena punto -> oggetto -> classe, in verticale: le etichette stanno a fianco
  % degli archi e non toccano i nodi
  \node[iri] (p) at (0,0)      {\texttt{base:9805241}};
  \node[iri] (o) at (0,-1.55)  {\texttt{base:\_Capital\_number\_10\_}};
  \node[cls] (c) at (0,-3.10)  {\texttt{base:Capital}};

  \draw[arc] (p) -- (o)
    node[side, anchor=west, pos=0.5, xshift=2mm] {\texttt{base:Constitutes}};
  \draw[arc] (o) -- (c)
    node[side, anchor=west, pos=0.5, xshift=2mm] {\texttt{rdf:type}};

  % coordinate e colore: una tripla per letterale, ben distanziate perche' le
  % quattro etichette abbiano spazio lungo gli archi
  \node[lit] (x)   at (-5.7, 1.35) {\texttt{12.4}};
  \node[lit] (y)   at (-5.7, 0.45) {\texttt{3.1}};
  \node[lit] (z)   at (-5.7,-0.45) {\texttt{88.0}};
  \node[lit] (rgb) at (-5.7,-1.35) {\texttt{210} \texttt{198} \texttt{173}};

  \draw[arc] (p) -- (x)   node[side, pos=0.45, above] {\texttt{base:X}};
  \draw[arc] (p) -- (y)   node[side, pos=0.45, above] {\texttt{base:Y}};
  \draw[arc] (p) -- (z)   node[side, pos=0.45, below] {\texttt{base:Z}};
  \draw[arc] (p) -- (rgb) node[side, pos=0.45, below] {\texttt{base:R/G/B}};
\end{tikzpicture}
